\documentclass[a4paper, 12pt]{article}

\usepackage{utils}

\renewcommand*{\today}{29 January 2026}

\begin{document}

\hotbox{Analyse 4}{CM 4}{\today}

\section{Séries de fonctions}

Dans cette section, nous considérons la suite de fonctions $(f_n)$ définie sur un intervalle $I \subset \mathbb{R}$ à valeurs dans $\R$ (ou $\C$).

De plus on définit $\funcdef{\theta}{I}{\R}{x}{0}$ la fonction nulle.

\begin{definition}
    On appelle \textbf{série de fonctions} de terme général $f_n$, notée $\sum\limits_k f_k$, la suite des $(S_n)_{n \in \N}$, appelées \textbf{sommes partielles}, définie par
    $$
    S_n(x) = \sum_{k=k_0}^n f_k(x)
    $$
\end{definition}

\begin{definition}
    On dira que la série $\sum\limits_k f_k$ 
    \begin{itemize}[label=--]
        \item 
        \textbf{converge simplement} vers $S$ sur $I$ si $(S_n)$ converge simplement vers $S$ sur $I$.
        \item 
        \textbf{converge uniformément} vers $S$ sur $I$ si $(S_n)$ converge uniformément vers $S$ sur $I$.
    \end{itemize}
    où $S = \lim\limits_{n \to +\infty} S_n = \sum\limits_{k=k_0}^{+\infty} f_k(x)$.

    On appelle \textbf{reste} d'ordre $n$ de $\sum\limits_k f_k$ la quantité $\funcdef{R_n}{I}{\R}{x}{\sum_{k=n+1}^{+\infty} f_k(x)}$.
\end{definition}

\begin{exemple}
    Pour tout $n \in \N$, on définit $\funcdef{f_n}{\R}{\R}{x}{x^n}$.

    On a la somme partielle $S_n(x) = \sum\limits_{k=0}^n x^k = \dfrac{1-x^{n+1}}{1-x}$ si $x \neq 1$

    On peut distinguer plusieurs cas:

    \begin{itemize}[label=--]
        \item si $|x| < 1, \lim\limits_{n \to +\infty} S_n(x) = \frac{1}{1-x} =: S(x)$.
        \item si $x = 1$ $f_n(x) = 1$ et donc $S_n(1) = \sum\limits_{k=0}^n f_n(1) = n + 1 \to +\infty$ donc $S(1)$ n'existe pas.
        \item si $x = -1$ on a $f_n(x) = (-1)^n$ et donc \text{$S_n(-1) = \sum\limits_{k=0}^n f_k(-1) = \begin{cases}
            0 & \text{si } n \text{ est impair}\\
            1 & \text{si } n \text{ est pair}
        \end{cases}$} donc $S(-1)$ n'existe pas.
        \item si $|x| > 1$, $\lim\limits_{n \to +\infty} S_n(x) = +\infty$.
    \end{itemize}

    Ainsi $\sum\limits_{k}f_k$ converge simplement vers $\dfrac{1}{1 - x}$ sur $]-1, 1[$

    \vspace{1cm}
    
    \begin{hotwarn}
        Demander au prof au niveau de la rigueur
    \end{hotwarn}

    Étudions sa convergence uniforme sur $]-1, 1[$.

    Pour un $n \in \N$ fixé et $x \in ]-1, 1[$, on a
    $$
    |S_n(x) - S(x)| = \left|\dfrac{1 - x^{n+1}}{1 - x} - \dfrac{1}{1 - x}\right| = \left|\dfrac{-x^{n+1}}{1 - x}\right| = \dfrac{|x|^{n+1}}{|1 - x|} \underset{x \to 1^-}{\to} +\infty
    $$
    car $\lim\limits_{x \to 1^-} |x|^{n+1} = 1$ et $\lim\limits_{x \to 1^-} |1 - x| = 0^+$.

    d'où pour tout $n \in \N$, $\norm{S_n - S}_{\infty, ]-1, 1[} = +\infty \underset{n \to +\infty}{\not\to} 0$ \par donc
    $(S_n)$ ne converge pas uniformément vers $S$ sur $]-1, 1[$.

    \vspace{1cm}

    Étudions la convergence uniforme sur tout compact \text{$J := [a, b] \subset ]-1, 1[$} avec $-1 < a < b < 1$.

    Pour un $n \in \N$ fixé et $x \in [a, b]$, on a
    $$
    |S_n(x) - S(x)| = \left|\dfrac{1-x^{n+1}}{1-x} - \dfrac{1}{1-x}\right| = \left|\dfrac{-x^{n+1}}{1-x}\right| = \dfrac{|x|^{n+1}}{1 - x} \leq \dfrac{K^{n+1}}{1 - b} \underset{n \to +\infty}{\to} 0
    $$
    avec $K = \max(|a|, |b|) < 1$

    car
    $$
    |x| \leq K < 1 \implies 0 < |x|^{n+1} \leq K^{n+1} < 1
    $$
    et
    $$
    x < b \implies 1 - x > 1 - b > 0 \implies 0 < \dfrac{1}{1 - x} \leq \dfrac{1}{1 - b}
    $$
    et aussi, $\lim\limits_{n \to +\infty} K^{n+1} = 0$.

    Comme $|S_n(x) - S(x)|$ est majoré par une quantité indépendante de $x$ qui tend vers $0$ quand $n$ tend vers $+\infty$, on en déduit que
    $(S_n)$ converge uniformément vers $S$ sur $J$, c'est-à-dire sur tout compact de $]-1, 1[$.
\end{exemple}

\begin{proposition}{}{}
    On suppose que $\sum\limits_k f_k$ converge simplement vers $S$ sur $I$.

    Alors la série de fonction $\sum\limits_k f_k$ converge uniformément vers $S$ sur $I$ si et seulement si
    la suite des restes $(R_n)$ converge uniformément vers $\theta$ sur $I$.
\end{proposition}

\begin{demonstration}
    $(S_n)$ converge uniformément vers $S$ sur $I$ si et seulement si $\norm{S_n - S}_{\infty, I} \to 0$.

    Or,
    \begin{align*}
        \norm{S_n - S}_{\infty, I} \to 0 &\iff \norm{\sum_{k=0}^n f_k - \sum_{k=0}^{+\infty} f_k}_{\infty, I} \to 0 \\
        &\iff \norm{\sum_{k=n+1}^{+\infty} f_k}_{\infty, I} = \norm{R_n - \theta}_{\infty, I} \to 0
    \end{align*}
\end{demonstration}

\begin{exemple}
    Pour tout $n \in \N$, on définit $\funcdef{f_n}{[-1, 1]}{\R}{x}{\dfrac{x^2}{n^2}}$ et on étudie $\sum\limits_k f_k$.

    Pour $x \in [-1, 1]$ on a 
    $$
    \left|\dfrac{x^2}{n^2} - \theta(x)\right| = \left|\dfrac{x^2}{n^2}\right| < \dfrac{1}{n^2} \to 0
    $$
    c'est-à-dire, convergence uniforme de $f_n$ vers $\theta$ sur $[-1, 1]$.

    On a
    \begin{align*}
        R_n(x) - \theta(x) = R_n(x) &= \sum_{k=n+1}^{+\infty} f_k(x) \\
        &= \sum_{k=n+1}^{+\infty} \dfrac{x^2}{k^2} \\
        &= x^2 \sum_{k=n+1}^{+\infty} \dfrac{1}{k^2} \\
        &\leq \sum_{k=n+1}^{+\infty} \dfrac{1}{k^2}
    \end{align*}

    Or, $\sum\limits_{k=n+1}^{+\infty} \dfrac{1}{k^2}$ est le reste d'une série de Riemann, qui converge donc $\lim\limits_{n \to +\infty} \sum\limits_{k=n+1}^{+\infty} \dfrac{1}{k^2} = 0$.

    Donc $R_n$ converge uniformément vers $\theta$ sur $[-1, 1]$, et par conséquent, la série de fonctions $\sum\limits_k f_k$ converge uniformément sur $[-1, 1]$.
\end{exemple}

\end{document}